%% kapa: exam=1 round=37 year=2026 period=1 form=A subject=economics subject_ko=경제학원론 seq=25 number=65
%% source: Q-Net 감정평가사 기출문제 (cst00309) artlSeq=5258924 file=raw/1008/5258924/2026년 제37회 감정평가사 1차 1교시.pdf
%% conversion: kordoc 4.24.51 -> tools/gen1.py (auto); review=proofread-sonnet; flags=none
\begin{question}{65}
화폐수요이론에 관한 설명으로 옳은 것은?
\begin{choices}
  \item 화폐수요의 이자율 탄력성이 0이 될 때 유동성 함정이 나타난다.
  \item 케인즈의 거래적 동기의 화폐수요는 이자율과 관련이 있다.
  \item 화폐의 중립성에 의하면 통화정책으로 물가를 관리할 수 없다.
  \item 고전학파의 화폐수량설에서는 화폐의 유통속도는 일정하지 않다고 가정한다.
  \item 유동성선호이론에 의하면 명목이자율은 화폐 보유의 기회비용이다.
\end{choices}
\end{question}
\answer{5}
